\section{Proof of the four-generator inequality}\label{sec:proof}

All vector comparisons are coordinatewise. A lower ideal contains every
nonnegative integer vector below each of its points. A positive corner
means a minimal excluded point with all three coordinates positive.

\subsection{Apéry representatives, moments, and collisions}\label{sec:p1}

The Apéry set
\[
\operatorname{Ap}(S,m)=\{w\in S:w-m\notin S\}
\]
consists of the least semigroup element in each residue modulo \(m\).
It has \(m\) elements, every factorization of which avoids \(m\).
Minimality of the generators gives \(0,a_1,a_2,a_3\in\operatorname{Ap}(S,m)\),
with the generator factorizations unique. Its largest element is
\[
M=c+m-1:
\]
every \(w-m\) is at most the largest gap \(c-1\), while \(c-1+m\) is Apéry.

For each Apéry element choose the lexicographically least exponent vector
\(x\in\mathbb N^3\) with label \(\lambda(x)=a\cdot x=w\).
Let \(T\) be the chosen vectors. This is the preferred-factorization
construction of Zhai, with the initial-ideal interpretation of
Hellus--Rechenauer--Waldi \cite{Z,HRW}.

To prove lower closure, take \(y\le x\in T\) and \(v=x-y\).
If \(\lambda(y)-m\in S\), then \(\lambda(x)-m\in S\), a contradiction.
Thus \(\lambda(y)\) is Apéry. A lexicographically earlier factorization
\(y'\) would give \(y'+v<_{\rm lex}x\) of the same label. Consequently
\(y\in T\). In particular,
\[
|T|=m,\quad 0,e_1,e_2,e_3\in T,\quad
\lambda:T\longrightarrow\operatorname{Ap}(S,m)
\ \text{is bijective},
\]
and the labels on \(T\) are distinct modulo \(m\).

Let \(s=\sum_Tx\), \(w=a/A\), \(H=M/A=\max_Tw\cdot x\), and
\(R=\max_T|x|_1\). Every \(w_i\ge1\), so \(R\le H\).
Writing the least representative of residue \(r\) as \(r+k_rm\)
and counting its \(k_r\) gaps proves the classical genus formula
(also \cite[Lemma~1]{Z})
\[
g=\frac1m\sum_{\operatorname{Ap}}w-\frac{m-1}{2}.
\]
Since \(W_4=3c-4g\), substitution gives the exact identity
\begin{equation}
\boxed{mW_4=A D_0-m(m-1),\qquad D_0=3mH-4w\cdot s.}
\label{eq:p:1}
\end{equation}

A minimal excluded point, or \emph{corner}, is \(p\notin T\) with
\(p-e_i\in T\) whenever \(p_i>0\). Let \(\rho(p)\in T\) have its residue.
If both \(p_i,\rho(p)_i>0\), subtracting \(e_i\) gives distinct included
points of equal residue, impossible. Likewise two distinct corners
sharing a positive coordinate cannot have the same residue.
Hence a positive corner has representative zero, and there is at most one
positive corner. All other corners have a zero coordinate.
These are the support/collision consequences of \cite[Proposition~2.6]{HRW}.

\subsection{Coordinate lines and two centroid bounds}\label{sec:p2}

For each of the three coordinate directions partition \(T\) into its
nonempty coordinate lines. Use an \emph{indexed} line family: intersecting
geometric lines in different directions remain distinct.
A line of length \(L_\ell\) starts at coordinate zero and has top \(t_\ell\).
Arithmetic progression sums give
\begin{equation}
\sum_\ell L_\ell=3m,\qquad
\sum_\ell L_\ell t_\ell=4s.
\label{eq:p:2}
\end{equation}
In a fixed direction its varying-coordinate contribution is twice that
coordinate's moment; the other coordinates contribute their moments once.
Thus, for any positive weights and any allowance \(Q\ge\max_Tw\cdot x\),
\begin{equation}
D_Q:=3mQ-4w\cdot s
=\sum_\ell L_\ell(Q-w\cdot t_\ell)\ge0.
\label{eq:p:3}
\end{equation}
This is the dimension-three weighted downset baseline of \cite[Lemma~3]{Z}.

Define
\[
\rho_2(t)=\max\{|v|_1:v\in\mathbb N^3,\ |v|_1\le2,\ t+v\in T\},
\qquad U_2=\sum_\ell L_\ell\rho_2(t_\ell).
\]
Choose a maximizing destination for each line. The total destination mass
is \(3m\), and its moment is \(4s+\beta\), with \(\beta\ge0\)
coordinatewise and \(|\beta|_1=U_2\). Each destination has weight at most
\(Q\). For weights at least one this proves
\begin{equation}
D_Q\ge U_2.
\label{eq:p:4}
\end{equation}
Only ten offsets are possible: zero, three units, three doubled units,
and three sums of distinct units.

More generally, any nonnegative point masses \(\alpha_x\) on \(T\)
with total \(3m\) and residual
\(\beta=\sum_x\alpha_xx-4s\ge0\) give \(D_Q\ge|\beta|_1\).
Indeed
\(D_Q=\sum_x\alpha_x(Q-w\cdot x)+w\cdot\beta\).
This identity also proves the explicit exceptional witnesses below.

If \(q_i=\max_Tx_i>0\), \(q_*=\max_iq_i\), put
\begin{equation}
G=q_*\left(3m-4\sum_i s_i/q_i\right).
\label{eq:p:5}
\end{equation}
When \(C=3m-4\sum_i s_i/q_i\ge0\), assign mass \(4s_i/q_i\)
to \(q_ie_i\), and mass \(C\) to a longest-axis endpoint.
The total is \(3m\), with residual \(Cq_*e_j\ge0\), proving \(D_Q\ge G\).
When \(C<0\), \eqref{eq:p:3} gives \(D_Q\ge0>G\).
In all cases \(D_Q\ge\max(U_2,G)\).

\subsection{The final-window projection inequality}\label{sec:p3}

Let \(\pi_j\) delete coordinate \(j\), and put
\[
P(T)=\sum_j|\pi_jT|,\quad
Z=\{x\in T:M-\lambda(x)<m\},\quad K=\operatorname{Max}(T).
\]
The set \(Z\) is nonempty. Its points are maximal, since an included
successor would increase their labels by \(a_j>m\) beyond \(M\).
For a nonempty point set \(X\), define
\[
E(X)=\sum_X|x|_1-\sum_j\max_Xx_j,\qquad
\Phi(T,X)=P(T)-3|X|-E(X).
\]
Here \(E(X)\ge0\). Adding a point increases each coordinate sum at least
as much as its coordinate maximum, so \(E(X)\le E(Y)\) for \(X\subseteq Y\).
Consequently \(\Phi(T,Z)\ge\Phi(T,K)\).

We prove
\begin{equation}
\boxed{mW_4\ge m\Phi(T,Z)+E(Z).}
\label{eq:p:6}
\end{equation}
Use the unnormalized weights \(a\) in the line formula, and write
\(M-\lambda(t_\ell)=m d_\ell+\beta_\ell\), \(0\le\beta_\ell<m\).
For a direction-\(j\) line of length \(L\), define
\[
R_j(L)=\sum_{r=0}^{L-1}[ra_j]_m,\quad
C_\ell=\#\{r<L:\beta_\ell+[ra_j]_m\ge m\},\quad
\sigma_\ell=mL d_\ell+mC_\ell-R_j(L).
\]
Reading the line backward gives
\[
L(M-\lambda(t_\ell))
=\sum_{x\in\ell}[M-\lambda(x)]_m+\sigma_\ell.
\]
Each direction partitions all \(m\) residues. Summing and using \eqref{eq:p:1} gives
\(mW_4=\binom m2+\sum_\ell\sigma_\ell\).

A deep line, whose top is outside \(Z\), has \(d_\ell\ge1\).
Its positive axis-prefix residues are distinct nonzero residues on \(T\);
therefore \(R_j(L)\le m(L-1)-\binom L2\) and \(\sigma_\ell\ge m\).
There are \(P(T)-3|Z|\) deep lines, because every \(z\in Z\) tops exactly
one line in each direction.

On a final line \(d_\ell=0\), so \(\sigma_\ell\ge-R_j(L)\).
Across final lines each positive axis-prefix point appears once, with
exactly \(E(Z)\) repeated occurrences: its multiplicity in direction \(j\)
is \(\#\{z\in Z:z_j\ge r\}\).
Distinct axis-prefix points have distinct nonzero residues.
The sum \(\binom m2\) pays for one copy of each used residue;
each repeated copy costs at most \(m-1\).
Thus
\[
mW_4\ge m(P(T)-3|Z|)-(m-1)E(Z),
\]
which is \eqref{eq:p:6}.
In particular \(\Phi(T,K)\ge0\) suffices to prove Wilf.

\subsection{Phase and rectangular thickening}\label{sec:p4}

Normalize by attained height: \(u_i=w_i/H\), \(v=\min_i u_i\),
\(U=\max_i u_i\), \(\sigma=\sum_i u_i\),
\(\mu_i=\mathbb E_T(u_iX_i)\), and
\(\kappa=3-4\sum_i\mu_i=D_0/(mH)\).
Units give \(u_i\le1\). The mean slack of direction \(j\) line tops is
defined as follows. For \(x\in T\), let \(t_j(x)\) be the top of the
direction-\(j\) line through \(x\), and put
\[
\delta_j(x)=1-u\cdot t_j(x),\qquad
\bar\delta_j=\frac1m\sum_{x\in T}\delta_j(x).
\]
Thus \(\bar\delta_j\) is the length-weighted mean over direction-\(j\)
lines, and the line moment identities give
\[
\bar\delta_j=(1+\kappa)/4-\mu_j,\qquad
\sum_j\bar\delta_j=\kappa\ge0.
\]

Choose coordinate \(k\) of weight \(v\).
Every interval of length \(h\) has sawtooth integral
\[
\int_I[r]_q\,dr\ge\frac{qh^2}{2(h+q)}.
\]
To prove this, write \(h=nq+r_0\): \(n\) full periods contribute
\(nq^2/2\), the remaining arc at least \(r_0^2/2\);
Cauchy--Schwarz gives \((nq^2+r_0^2)/2\ge h^2/(2(n+1))\),
and \(n+1\le(h+q)/q\).

On a coordinate-\(k\) line of length \(L\), write \(\delta_j(t)\) for the
slack at its point with \(k\)-coordinate \(t\). For \(j\ne k\), these
slacks have residues \([\theta-tv]_{u_j}\).
Nonnegative slack and
\([\theta-tv+r]_{u_j}\le\delta_j(t)+r\), \(0\le r\le v\),
bound the integral over the concatenated intervals of total length \(Lv\).
Since \(Lv\le1+v\), division gives
\[
\operatorname{mean}_{\rm line}\delta_j
\ge\frac{u_jvL}{2(1+v+u_j)}-\frac v2.
\]
Length-weighted averaging uses \(v\sum L^2/m=2\mu_k+v\), hence
\[
2\bar\delta_j(1+v+u_j)\ge2u_j\mu_k-v(1+v).
\]
Sum over \(j\ne k\), substitute \(\mu_k=(1+\kappa)/4-\bar\delta_k\),
and use \(U\le\sigma-2v\).
The remaining slack coefficient is the middle weight minus \(1+v\),
nonpositive because that weight is at most one. Discarding it gives
\begin{equation}
\boxed{\sigma(1-3\kappa)\le4\kappa+5v-3\kappa v+4v^2.}
\label{eq:p:7}
\end{equation}
If \(H>3\) and \(D_0<m\), then \(0\le\kappa<v=1/H<1/3\).
The right side divided by \(1-3\kappa\) is strictly increasing in \(\kappa\),
with derivative numerator \(4+12v+12v^2>0\).
Writing \(B=\sum_iw_i\) proves the necessary failure cutoff
\begin{equation}
\boxed{B<9+28/(H-3).}
\label{eq:p:8}
\end{equation}

Thicken \(T\) into equal-volume half-open boxes
\[
K_T=\bigcup_{x\in T}\prod_i
\left[\frac{w_ix_i}{H+B},\frac{w_i(x_i+1)}{H+B}\right).
\]
It is a positive-volume downset in the unit simplex.
Its mean coordinate sum is
\((w\cdot s/m+B/2)/(H+B)\), giving
\begin{equation}
\boxed{\kappa_c(K_T):=3-4\mathbb E_{K_T}|Y|_1
=\frac{D_0/m+B}{H+B}.}
\label{eq:p:9}
\end{equation}
The complement is the union of upper orthants with vertices
\((w_ip_i/(H+B))_i\), for the discrete minimal exclusions \(p\).
Positive diagonal scaling preserves their supports and minimality.
The same calculation with an allowance \(Q\) in place of \(H\) replaces
\(D_0\) by \(D_Q\) in the identity.
A continuous lower bound \(\kappa_c\ge\gamma\) therefore gives
\begin{equation}
D_0/m\ge\gamma H-(1-\gamma)B,\qquad
D_0<m\ \Longrightarrow\ \gamma H<1+(1-\gamma)B.
\label{eq:p:10}
\end{equation}

\begin{lemma}[Height cutoff]\label{lemma:height-cutoff}
Suppose \(H>3\), \(D_0<m\), and
\[
\frac{D_0}{m}\ge\alpha H-\beta B,\qquad
\alpha>0,\quad 0\le\beta\le3.
\]
Then
\begin{equation}
\boxed{H<C:=3+\frac{1+9\beta}{\alpha}}.
\label{eq:height-cutoff}
\end{equation}
\end{lemma}
\begin{proof}
The assumed moment bound gives \(\alpha H<1+\beta B\).
Substitute \eqref{eq:p:8} and multiply by the positive quantity \(H-3\):
\[
\alpha H^2-(3\alpha+1+9\beta)H+3-\beta
=\alpha H(H-C)+3-\beta<0.
\]
If \(H\ge C\), the left side is nonnegative because \(\alpha>0\)
and \(3-\beta\ge0\), a contradiction.
\end{proof}

\begin{table}[htbp]
\centering
\begin{tabular}{@{}lccc@{}}
\toprule
Branch & \(\alpha\) & \(\beta\) & Bound \\
\midrule
No full-support corner & \(1/3\) & \(2/3\) & \(H<24\) \\
Short corner, high height & \(1/3\) & \(4/3\) & \(H<42\) \\
Residual high height & \(5/42\) & \(37/42\) & \(H<78\) \\
\bottomrule
\end{tabular}
\caption{Parameters for the common height-cutoff argument. The interval
certificates cover the corresponding closed supersets.}
\label{tab:height-cutoffs}
\end{table}

\subsection{Structural geometry and an exact interval method}\label{sec:p5}

Theorem~\ref{theorem:a:1.1}, proved in Section~\ref{sec:analytic-a},
proves that every finite no-full-corner lower ideal \(U\) is a disjoint
central anchored box \([0,z]\) and at most three outward horns.
Horn \(i\) consists of rectangles at levels \(t>z_i\), with transverse
caps at most \(z_j,z_k\), nonincreasing with \(t\).
The proof constructs a chordal compatibility graph, proves its needed
clique-tree property, and uses the median of three coordinate maximizers.

For a unique positive corner \(p\), deleting its generator from the
upper-ideal complement gives a finite no-full-corner \(U\) with
\begin{equation}
T=U\setminus(p+\mathbb N^3).
\label{eq:p:11}
\end{equation}
Pure-axis exclusions remain, so \(U\) is finite and has the same axes.
Clipping the center and horns preserves their disjointness.

We use the following exact upper-bound method in the three interval checks.
For these checks, permute coordinates and weights together so that
\(w=(1,b,d)\), with \(1\le b\le d\).
At scale \(q=4096\), integer endpoints \(L=(B_0,C_0,Q_0)\),
\(V=(B_1,C_1,Q_1)\) enclose a closed real parameter box for \((b,d,Q)\).
Put
\[
\ell=(q,B_0,C_0),\quad u=(q,B_1,C_1),\quad
N_i=\lfloor Q_1/\ell_i\rfloor,\quad
\psi(x)=4u\cdot x+q-3Q_0.
\]
Every relevant retained point satisfies \(\ell\cdot x\le Q_1\).
Every axis of \(U\) survives clipping, so \(U\subseteq\prod_i[0,N_i]\).
Pointwise monotonicity gives, for every real parameter in the box,
\begin{equation}
q\bigl(m-(3mQ-4w\cdot s)\bigr)\le\sum_T\psi(x).
\label{eq:p:12}
\end{equation}

For an inclusive box \([a,d]\), its count is \(\prod_i(d_i-a_i+1)\)
and twice its objective moment is that count times \(\sum_i u_i(a_i+d_i)\).
For a clipped box subtract the upper box \([\max(a,p),d]\).
The retained feasibility maximum is
\begin{equation}
\max_{\substack{i\\a_i\le\min(d_i,p_i-1)}}
\left(\ell_i\min(d_i,p_i-1)+\sum_{j\ne i}\ell_jd_j\right).
\label{eq:p:13}
\end{equation}
This maximizes over the parts with some coordinate below \(p\).
An independent formula partitions by the \emph{first} such coordinate,
summing disjoint boxes with earlier coordinates at least their \(p_j\).
Their maximum heights and summed moments give the same statistics.
Without clipping use the full box and its upper-corner height.

Let \(C_i(t,r,s)\) be the exact section score, and call it feasible
when its retained maximum height is at most \(Q_1\).
For a possibly empty nested horn use
Lemma~\ref{lemma:horn-recurrence}, terminating beyond its axis cap.
Enumerate every feasible retained central box and add
its score to \(\sum_iF_i(z_i+1;z_j,z_k)\).
The structural theorem then proves an upper bound \(V_p(L,V)\) on \eqref{eq:p:12}.
Independent horns may enlarge the class; they never decrease its upper bound.

Order tightening replaces endpoints by
\begin{equation}
L^\sharp=(B_0,\max(B_0,C_0),\max(B_0,C_0,Q_0)),\quad
V^\sharp=(\min(B_1,C_1,Q_1),\min(C_1,Q_1),Q_1).
\label{eq:p:15}
\end{equation}
Its ordered intersection is unchanged; inversion proves emptiness.
Splitting a tightened box at an interior integer wall yields two closed
children sharing that wall and covering the parent.
A finite tree reaching every node exactly once and ending only at checked
empty, analytic, or exact-DP leaves covers all real root parameters.


\subsection{Small multiplicity}\label{sec:p6}

The two published computer-assisted results \cite[Theorem~4.3]{B} and
\cite[Proposition~6.9]{KS} settle all numerical semigroups of multiplicity at
most 18 and exactly 19, respectively, without an embedding-dimension
restriction.

For \(20\le m\le29\) we need an analytic counterexample reduction.
Theorem~\ref{theorem:b:1.1}, proved in Section~\ref{sec:analytic-b},
proves Wilf whenever \(T=\{x:a\cdot x\le M\}\), including an analytic
residue obstruction to its formal six-column case.

If the ideal is not full weighted, choose an omitted point of weight
at most \(M\) of least degree. It is a minimal excluded point \(p\)
with \(a\cdot p\le M\). Its direction-\(i\) predecessor line, for \(p_i>0\),
has top \(p-e_i\) and length exactly \(p_i\).
Their contribution to the unnormalized line slack is
\[
\sum_{i:p_i>0}p_i(M-a\cdot p+a_i)
=M+(|p|_1-1)(M-a\cdot p)\ge M.
\]
All other slacks are nonnegative. Therefore \(W_4<0\), which implies
the integer bound \(W_4\le-1\), forces
\begin{equation}
\boxed{M\le3mM-4a\cdot s\le m(m-2).}
\label{eq:p:16}
\end{equation}
Every generator \(a_i\) is Apéry and at most \(M\).
After sorting, any counterexample thus belongs to the \emph{inclusive} box
\[
m<a_1<a_2<a_3\le B_m:=m(m-2).
\]

\begin{finitelemma}[Generator-box]\label{finite:generator-box}
For each \(20\le m\le29\), every increasing
triple in this box with \(\gcd(m,a_1,a_2,a_3)=1\) and all four generators
minimal has either largest Apéry element \(M>B_m\), or \(W_4\ge0\).

\end{finitelemma}

Every raw triple is classified by successive exact tests:
\(a_1\in\langle m\rangle\), \(a_2\in\langle m,a_1\rangle\),
\(a_3\in\langle m,a_1,a_2\rangle\), then a nontrivial gcd, or validity.
Larger generators cannot make an earlier one redundant.
One path computes exact residue shortest distances \(d_r\), adjoining
a generator by two forward passes around each residue cycle.
A shortest extension uses fewer than the cycle length many new copies;
two passes propagate every possible starting residue around the cycle.
The second path uses increasing ordinary membership:
\(I_j(n)=I_{j-1}(n)\vee(n\ge a_j\ \&\ I_j(n-a_j))\).
It tests the complete window \([B_m-m+1,B_m]\), which is full exactly
when \(M\le B_m\). A full window and closure under \(+m\) prove every
larger integer is included, so gaps and conductor are fully determined.

The generator-box finite lemma contradicts \eqref{eq:p:16} for any negative example.
Thus all multiplicities through 29 are settled.

\subsection{Large multiplicity: exhaustive routing and the three-plane case}\label{sec:p7}

Now \(m\ge30\). If there is no positive corner use Section~\ref{sec:p8}.
Otherwise let \(p\) be the unique positive corner.
If \(|p|_1=3\), then \(p=(1,1,1)\) and the following analytic argument applies.
Degree four gives a permutation of \((2,1,1)\), handled in Section~\ref{sec:p9}.
For \(|p|_1\ge5\), degree \(R\le6\) uses Section~\ref{sec:p10}; otherwise \(R\ge7\)
uses Section~\ref{sec:p11}. These are disjoint exhaustive alternatives.
Minimality gives \(|p|_1-1\le R\); consequently Section~\ref{sec:p10} has
\(5\le|p|_1\le7\), while Section~\ref{sec:p11} has \(H\ge R\ge7\).
All coordinate permutations are applied simultaneously to weights.

Suppose \(p=(1,1,1)\). Every point of \(T\) lies in a coordinate plane.
A mixed corner \(q\) in the \(ij\)-plane has representative \(\gamma e_k\),
\(1\le\gamma\le h_k=\max_Tx_k\). The included point
\(q-e_i-e_j\) has residue \((\gamma+1)a_k\), since \(a\cdot p\equiv0\).
If \(\gamma<h_k\) it collides with \((\gamma+1)e_k\).
Thus all mixed corners in that plane have representative \(h_ke_k\);
the collision rule permits at most one.

A geometric count proves \(\Phi(T,K)\ge0\) for \(m\ge30\).
Write \(h=h_1+h_2+h_3\ge3\).
Each plane is its full rectangle or that rectangle with one upper-right
quadrant deleted; it contains \((1,1)\).
Its frontier has one of the following forms, allowing transpose:

\begin{center}
\small
\begin{tabular}{@{}lllll@{}}
\toprule
Type & Mixed maxima & Mixed point count \(c_{ij}\) & Correction \(\beta_{ij}\) & Axis charge \(\nu_{ij}\) \\
\midrule
Rectangle & \((h_i,h_j)\) & \(h_ih_j\) & 0 & 0 \\
Split, \(1\le b<h_i,\ 1\le d<h_j\) & \((b,h_j),(h_i,d)\) & \(bh_j+dh_i-bd\) & \(b+d\) & 0 \\
Side cap at \(i\), \(1\le b<h_i\) & \((b,h_j)\) & \(bh_j\) & \(b\) & \(h_i\) \\
\bottomrule
\end{tabular}
\end{center}

Let \(s_f\) count split planes, \(I\) the globally maximal axis directions,
\(J=\sum_{i\in I}h_i\), \(\Beta=\sum\beta_{ij}\),
\(C=\sum(c_{ij}-\beta_{ij})\), and \(N=\sum\nu_{ij}\).
Every split plane requires both incident axis caps to be at least two;
checking \(s_f=0,1,2,3\) gives \(h\ge s_f+3\).
Every mixed planar maximum is globally maximal. An axis top is globally
maximal exactly when both incident planes side-cap that axis.
Each side-cap plane can cap only one axis, while a split plane caps none.
Consequently \(2|I|\le3-s_f\), so \(s_f+|I|\le3\) and
\(k=|K|=3+s_f+|I|\le6\).
For a split,
\(c_{ij}-\beta_{ij}
=bd+b(h_j-d-1)+d(h_i-b-1)\ge b+d-1\).
For a side cap \(c_{ij}-\beta_{ij}\ge0\) and \(\beta_{ij}\le\nu_{ij}\).
Hence
\[
\Beta\le C+s_f+N,\qquad N\ge2J.
\]
Direct origin/axis/plane counting yields
\[
m=1+h+\Beta+C,\quad P(T)=m+h+2,\quad E(K)=h+\Beta-N+J.
\]
Set \(Q_*=h+C+N-J\). Then \(\Phi(T,K)=Q_*+3-3k\).
If this were negative, integrality would give \(Q_*\le3k-4\le14\), and
\[
m=1+Q_*+\Beta-N+J
\le1+2Q_*-h+s_f-N+2J
\le2Q_*-2\le26.
\]
This contradicts \(m\ge30\). Therefore \eqref{eq:p:6} proves \(W_4\ge0\).

\subsection{No full-support corner}\label{sec:p8}

Theorem~\ref{theorem:c:5.2}, proved in Section~\ref{sec:analytic-c},
proves \(\kappa_c\ge1/3\) for every positive-volume finite-step
central-box/horn set in the unit simplex.
That section proves that degree at most four in this discrete class
implies cardinality at most 29, by eleven explicit central-cap bounds.

If \(D_0<m-1\) at \(m\ge30\), then \(H\ge5\), since \(R\le H<5\)
would force degree at most four. Sort weights to \((1,b,d)\), so
\(1\le b\le d\le H\). The continuous bound and \eqref{eq:p:10}
give the first row of Table~\ref{tab:height-cutoffs}, so
Lemma~\ref{lemma:height-cutoff} gives \(H<24\).
For any allowance \(Q\), the continuous theorem also gives
\(D_Q/m\ge(Q-2B)/3\).

\begin{finitelemma}[No-corner interval]\label{finite:no-corner-interval}
For all real
\(1\le b\le d\le Q\), \(5\le Q\le24\), and every nonempty finite
no-full-corner ideal of attained height at most \(Q\),
\[
m-D_Q\le1.
\]

\end{finitelemma}

This is stronger than the genuine branch domain: no units, residue
labels, attained allowance, or cardinality restriction is required.
Use the method in Section~\ref{sec:p5}, with no clipping.
A DP leaf requires its exact bound at most \(q\).
An analytic leaf uses
\(Q_0\ge2(q+B_1+C_1)+3q\), which implies \(D_Q\ge m\).
The interval certificate covers the closed root from height 5 to 24.
The finite lemma contradicts the hypothetical failure.
Thus \(D_0\ge m-1\) throughout this branch.

\subsection{The short corner}\label{sec:p9}

Orient \(p=(2,1,1)\), without prescribing the least weight's direction.
For any degree allowance \(r\ge3\),
\begin{equation}
|T|\le2r^2-r+3.
\label{eq:p:17}
\end{equation}
Indeed the degree-\(r\) simplex minus the \(p\)-orthant has \(2r^2+2\)
points. For each \(1\le j\le r-1\), the excluded point \((1,j,r-j)\)
cannot dominate the unique full-support corner.
It must dominate a planar corner; at least one of
\((0,j,r-j),(1,0,r-j),(1,j,0)\) is therefore omitted.
These \(r-1\) triples are pairwise disjoint, forcing \eqref{eq:p:17}.
In particular \(H<6\) gives \(R\le5\) and \(30\le m\le48\).
The endpoint \(H=6\) belongs to the high-height case.

\subsubsection{High height}

For a translated planar lower ideal of offset \(h\) in height \(H\),
the dimension-two line identity gives mean weight at most \((2H+h)/3\).
Partition \(T\) into \(y=0\); \(y\ge1,z=0\);
\(x=0,y,z\ge1\); and \(x=1,y,z\ge1\).
These are translated planar lower ideals, with offsets respectively
\(0,w_2,w_2+w_3,B\). Thus \(D_0/m\ge(H-4B)/3\).
If \(H\ge6\) and \(D_0<m-1\), Lemma~\ref{lemma:height-cutoff}
with the second row of Table~\ref{tab:height-cutoffs} gives \(H<42\).
We retain the closed height-42 superset.

\begin{finitelemma}[Short-corner high]\label{finite:short-corner-high}
For all
\(1\le b\le d\le Q\), \(6\le Q\le42\), each coordinate permutation of
\(p=(2,1,1)\), and every nonempty finite no-full-corner \(U\),
the clipped \(T=U\setminus(p+\mathbb N^3)\) of height at most \(Q\)
satisfies \(m-D_Q\le1\).

\end{finitelemma}

Apply Section~\ref{sec:p5}'s retained-height feasibility and clipped scores to all
three corners. A DP leaf requires every corner bound at most \(q\).
The planar rule \(Q_0\ge4(q+B_1+C_1)+3q\) is also sound, including equality.
The complete closed tree establishes the finite lemma, contradicting failure.

\subsubsection{Low height: exact profiles and two inline exceptions}

A mixed \(xy\) corner has representative \(\gamma e_3\ne0\).
If its \(x\)-coordinate is at least two, subtracting \(2e_1+e_2\)
gives an included point with residue \((\gamma+1)a_3\).
It forces \(\gamma=h_3\); at most one such corner exists.
The antichain permits at most one additional corner at \(x=1\).
The \(xz\) plane is identical. A \(yz\) corner yields residue
\((\gamma+2)a_1\) on subtracting \(e_2+e_3\), forcing
\(\gamma\in\{h_1-1,h_1\}\).
Hence each plane has at most two mixed corners, with the refined incident
plane restriction, and contains respectively \((2,1),(2,1),(1,1)\).

A degree-five plane profile is \(f=(f_0,\ldots,f_5)\),
\(0\le f_i\le6-i\), nonincreasing, encoding \(\{(i,j):j<f_i\}\).
Let \(\ell(f)\) count positive columns and
\(\Delta(f)=\{i:1\le i<\ell(f),\ f_i<f_{i-1}\}\).
These, and only these, are mixed corners; a drop to zero is a pure-axis corner.

For \(xy,xz\), require \(f_2\ge2\), \(|\Delta(f)|\le2\), and at most
one drop at index at least two. For \(yz\), require \(h_1\ge2\)
and \(|\Delta(h)|\le2\).
Use all compatible ordered triples \(f,g,h\), with
\[
\ell(f)=\ell(g),\quad f_0=\ell(h),\quad g_0=h_0,
\]
and reconstruct
\begin{equation}
T=\{(x,y,z):y<f_x,\ z<g_x,\ z<h_y,\ (x,y,z)\not\ge(2,1,1)\}.
\label{eq:p:18}
\end{equation}
Keep exactly those of degree at most five and \(30\le|T|\le48\).
Matching axes recover each plane section exactly.
The required plane entries include all three predecessors of \(p\).
Any extra reconstructed point in a genuine ideal would dominate a planar
minimal exclusion, violating its corresponding pair test.
Thus every genuine low-height ideal is recovered uniquely.
Profiles alone do not guarantee total degree: the literal degree check is essential.

\begin{finitelemma}[Short-corner low]\label{finite:short-corner-low}
Every such triple satisfies
\(U_2\ge m-1\), except possibly the two keys
\[
\begin{aligned}
E_0&=((4,3,3,2,0,0),(6,3,3,2,0,0),(6,4,2,2,0,0)),\\
E_1&=((6,3,3,2,0,0),(4,3,3,2,0,0),(4,4,2,2,1,1)).
\end{aligned}
\]

\end{finitelemma}

For \(E_0\), direct reconstruction gives \(m=30\), \(s=(27,25,34)\).

The included points \((0,0,5),(2,2,0),(3,1,0)\) with masses \(38,48,4\)
have total mass 90 and residual
\((108,100,190)-4s=(0,0,54)\).
Section~\ref{sec:p2} therefore gives \(D_Q\ge54\ge29\) for all admissible real weights.
The second ideal and witness are obtained by interchanging \(y,z\),
including transposition of the \(yz\) profile.

Together with \eqref{eq:p:4}, the finite lemma proves \(D_0\ge m-1\) in low height.
This completes the short-corner branch.

\subsection{Residual degree at most six}\label{sec:p10}

Let \(5\le P=|p|_1\le7\), \(R\le6\).
In an \(ij\)-plane, distinct mixed corners have distinct representatives
\(\gamma e_k\ne0\). At most \(p_i-1\) have coordinate \(q_i<p_i\),
and at most \(p_j-1\) have \(q_j<p_j\), by antichain distinct coordinates.
For the remaining corners, subtract \(p_ie_i+p_je_j\).
The included result has residue \((\gamma+p_k)a_k\);
collision with the \(k\)-axis forces \(\gamma>h_k-p_k\).
At most \(p_k\) distinct representatives remain. Each plane thus has at most
\begin{equation}
(p_i-1)+(p_j-1)+p_k=P-2
\label{eq:p:19}
\end{equation}
mixed corners.

Define \(F_i=\{x\in T:x+e_i\notin T\}\), \(n_k=1+\max_Tx_k\).
For every residue-bijective lower ideal,
\begin{equation}
|F_i\cap F_j|\le2n_k.
\label{eq:p:20}
\end{equation}
To prove it, let \(C_{ij}\) contain excluded points with both positive
\(i,j\) coordinates and both respective predecessors included.
In each nonempty fixed-\(k\) slice, a planar ideal with \(h\) maxima has
\(h-1\) mixed corners, by its constant positive-height blocks.
All \(n_k\) slices are nonempty, hence
\(|C_{ij}|=|F_i\cap F_j|-n_k\).
Subtracting an included predecessor forces the representative of
each \(q\in C_{ij}\) onto the \(k\)-axis, and distinct \(q\) have distinct
representatives by comparing their \(i\)-predecessors.
This injects \(C_{ij}\) into the \(n_k\) axis points, proving \eqref{eq:p:20}.
Slice corners need not be globally minimal; only these two predecessors are used.

Sort coordinates of \(p\), simultaneously permuting weights.
The nine corners are
\[
(1,1,3),(1,2,2),(1,1,4),(1,2,3),(2,2,2),
(1,1,5),(1,2,4),(1,3,3),(2,2,3).
\]
Use all nonincreasing seven-entry profiles \(0\le f_i\le7-i\).
In an \(ij\)-plane require \(f_{p_i}\ge p_j+1\) and at most \(P-2\)
positive height drops. Impose the three shared-axis compatibilities
as in \eqref{eq:p:18}, and reconstruct with its own \(p\)-orthant deleted.
Keep precisely degree at most six, cardinality at least 30,
and the three bounds \eqref{eq:p:20}. No maxima, realizability, modular, or erosion filter is used.
The same plane-recovery and excluded-corner argument as for \eqref{eq:p:18}
proves complete, unique coverage of genuine oriented ideals.

\begin{finitelemma}[Degree-six]\label{finite:degree-six}
On this complete geometric family,
\begin{equation}
\max(U_2,G)\ge m.
\label{eq:p:21}
\end{equation}

\end{finitelemma}

The local integer predicate and the denominator-cleared axis predicate
are exact. With \(d=q_1q_2q_3>0\), the latter is
\[
q_*\left(3md-4\sum_i s_i\prod_{j\ne i}q_j\right)\ge md.
\]
Section~\ref{sec:p2} and \eqref{eq:p:21} give \(D_0\ge m\) for this branch.

\subsection{Higher degree: a finite strip and a real interval bound}\label{sec:p11}

Assume \(R\ge7\), \(|p|_1\ge5\). We prove \(D_0\ge m-29/10\).
The first finite premise supplies a continuous inequality; that inequality
restricts a hypothetical failure to a compact real parameter region.
The second finite premise excludes the entire region.

\subsubsection{The three-allowance strip}

\begin{finitelemma}[Strip]\label{finite:strip}
For \(r\in\{18,19,20\}\), every finite lower
ideal \(T'\subseteq\mathbb N^3\) of degree at most \(r\), with at most
one positive minimal exclusion, satisfies
\begin{equation}
3r|T'|-4\sum_{x\in T'}|x|_1\ge4|T'|.
\label{eq:p:22}
\end{equation}
Empty ideals and missing coordinate units are allowed.

\end{finitelemma}

The empty ideal satisfies the inequality directly. For nonempty ideals,
the following finite reduction is exhaustive. If there is a positive minimal
corner \(p'\), its predecessors imply \(|p'|_1\le r+1\).
Remove that exclusion generator to obtain a finite no-corner completion:
pure-axis exclusions remain, and deleting a positive orthant leaves every
axis point unchanged. Thus its axis caps are at most \(r\).
If there is no positive corner, take \(p'=(1,1,r-1)\); its orthant is
inactive on the degree-\(r\) simplex.
By symmetry enumerate exactly
\[
1\le p'_1\le p'_2\le p'_3,\qquad |p'|_1\le r+1.
\]
No minimality filter is imposed on this larger family.

Apply the decomposition of Section~\ref{sec:analytic-a}, clipped by \(p'\), to the score
\(f_r(x)=4|x|_1+4-3r\). All center and transverse caps range from zero
to \(r\). A section or center is feasible when its \textbf{retained}
maximum degree is at most \(r\), not when its removed upper corner is.
Apply Lemma~\ref{lemma:horn-recurrence} with the section sum of this
score, retained-degree feasibility, and termination beyond the degree
allowance.
The maximum of center score plus its three horn bounds is an upper bound
for \(\sum_{T'}f_r=4|T'|-D_r(T')\). Every configuration must have this
bound at most zero. This is a search over corner/allowance configurations
and exact dynamic-programming states.
Singletons, zero caps, and inactive corners are included.

\subsubsection{Transfer to the continuous gap}

Let \(K\) be a positive-volume downset in the unit simplex whose complement
is a finite union of upper orthants with at most one positive minimal vertex.
Put \(h=1/21\), and for \(z\in[0,h)^3\) define
\[
T_z=\{y\in\mathbb N^3:z+hy\in K\},\quad
r_z=\left\lfloor(1-|z|_1)/h\right\rfloor,\quad
\epsilon_z=1-|z|_1-hr_z.
\]
For \(z\ne0\), \(0<|z|_1<3h\), so
\(r_z\in\{18,19,20\}\) and \(0\le\epsilon_z<h\).
The zero translation is a null set and is omitted throughout the integral;
no allowance-21 assertion is needed.
Each excluded vertex \(g\) induces the discrete vertex
\[
g'_i=\max(0,\lceil(g_i-z_i)/h\rceil).
\]
A zero coordinate stays zero. Consequently at most one induced minimal
vertex is positive, even after redundant vertices are removed.
For almost every translation boundary conventions agree; the exceptional
translations lie in finitely many coordinate hyperplanes modulo \(h\).
The simplex condition gives \(|y|_1\le r_z\) for every \(y\in T_z\).
The hypotheses of \eqref{eq:p:22} therefore hold for \(T_z\), including when
units are absent. Expansion gives
\[
\sum_{y\in T_z}(3-4|z+hy|_1)
 =hD_{r_z}(T_z)+(3\epsilon_z-|z|_1)|T_z|
 \ge(4h-|z|_1)|T_z|.
\]
Integrating the partition \(x=z+hy\), whose Jacobian is one, gives
\[
3-4\mathbb E_K|X|_1
\ge4h-\sum_i\mathbb E_K(X_i\bmod h).
\]
A coordinate fiber of a downset is an interval \([0,L]\) up to endpoints.
Write \(L=qh+t\), \(0\le t<h\). Since
\[
\int_0^L(x\bmod h)\,dx=(qh^2+t^2)/2\le hL/2,
\]
Fubini implies each remainder average is at most \(h/2\). Hence
\begin{equation}
\boxed{\kappa_c(K)\ge5/42.}
\label{eq:p:24}
\end{equation}
The averaged remainder estimate is essential; a pointwise estimate alone
does not give this constant. The rectangular thickening from Section~\ref{sec:p4}
preserves excluded supports and meets every hypothesis above.

\subsubsection{Necessary region for a failure}

Sort the normalized weights to \(w=(1,b,d)\), and write \(B=1+b+d\).
The predecessors of the positive corner contain the units, so
\(1\le b\le d\le H\), and \(H\ge R\ge7\).
If \(D_0<m-29/10<m\), Section~\ref{sec:p4} and \eqref{eq:p:24} imply
\[
B<9+\frac{28}{H-3},\qquad
\frac{D_0}{m}\ge\frac{5H-37B}{42},\qquad 5H<42+37B.
\]
Lemma~\ref{lemma:height-cutoff}, with the final row of
Table~\ref{tab:height-cutoffs}, gives \(H<78\).
Also \(p-e_1\in T\) gives \(w\cdot p\le H+1\).
It suffices to cover the closed necessary failure region
\begin{equation}
\mathcal F=\{(b,d,Q):
1\le b\le d\le Q,\ 7\le Q\le78,\
1+b+d\le9+28/(Q-3),\
5Q\le42+37(1+b+d)\}.
\label{eq:p:25}
\end{equation}

\begin{finitelemma}[High-height]\label{finite:high-height}
For every \((b,d,Q)\in\mathcal F\),
every positive integer \(p\) with \(|p|_1\ge5\) and
\((1,b,d)\cdot p\le Q+1\), and every nonempty finite no-corner lower ideal \(U\),
put \(T'=U\setminus(p+\mathbb N^3)\). If
\(\max_{T'}(1,b,d)\cdot x\le Q\), then
\begin{equation}
|T'|-\left(3Q|T'|-4(1,b,d)\cdot\sum_{T'}x\right)\le29/10.
\label{eq:p:26}
\end{equation}

\end{finitelemma}

No residue, cardinality, minimality, unit, or surface restriction is used
in this enlarged geometric family.
Use the following interval tightening. At scale \(q=4096\) write endpoints
\(L=(B_0,C_0,H_0)\), \(U=(B_1,C_1,H_1)\) for \((b,d,Q)\).
First intersect with order:
\[
L\leftarrow(B_0,\max(B_0,C_0),\max(B_0,C_0,H_0)),\quad
U\leftarrow(\min(B_1,C_1,H_1),\min(C_1,H_1),H_1).
\]
If not inverted, compute, in this order,
\begin{equation}
\begin{aligned}
S_{\max}&=9q+\left\lceil28q^2/(H_0-3q)\right\rceil,\\
B_1&\leftarrow\min(B_1,\lceil(S_{\max}-q)/2\rceil),\\
C_1&\leftarrow\min(C_1,S_{\max}-q-B_0),\\
S_{\mathrm{upper}}&=\min(q+B_1+C_1,S_{\max}),\\
H_1&\leftarrow\min(H_1,\lceil(42q+37S_{\mathrm{upper}})/5\rceil).
\end{aligned}
\label{eq:p:27}
\end{equation}
Perform at most twenty iterations, stopping on inversion or unchanged
endpoints. Every upper rounding is outward. The inequalities in \eqref{eq:p:25},
the lower bound \(qQ\ge H_0\), and \(b\le d\) prove that each step retains
every relevant real point. Convergence is not required for soundness.

For a nonempty tightened box use Section~\ref{sec:p5} with
\(\ell=(q,B_0,C_0)\), \(u=(q,B_1,C_1)\), and
\(\psi(x)=4u\cdot x+q-3H_0\). Enumerate the \textbf{full Cartesian} corner list
\[
p_i\ge1,\quad5\le|p|_1\le\lfloor H_1/q\rfloor+1,\quad
\ell\cdot p\le H_1+q.
\]
This contains every relevant corner, since \(\ell_i\ge q\).
An empty list is accepted only as vacuous.
For every listed corner the retained-center/three-horn bound \(V_p\)
must satisfy \(10V_p\le29q\).
The closed root is \(((q,q,7q),(78q,78q,78q))\).
Each child is derived from its tightened parent; each split is strictly
interior, and its two children share the split wall. Checked empty leaves
and checked dynamic-programming leaves exhaust all nodes.
Induction proves \eqref{eq:p:26} throughout the closed real parameter region.
At \(Q=H\), it contradicts the proposed failure and proves this branch.

\subsection{Assembly of the proof}\label{sec:p12}

\begin{proof}[Proof of Theorem~\ref{thm:main}]

The branches are exhaustive. Multiplicity at most 29 is settled by Section~\ref{sec:p6}.
For \(m\ge30\), Section~\ref{sec:p7} routes every ideal to exactly one of Sections~\ref{sec:p8}--\ref{sec:p11} or the analytic three-plane case. The three-plane case already gives
\(W_4\ge0\). On the remaining branches their conclusions give
\(D_0\ge m-29/10\).
This bound is positive, and \(A\ge m+1\). Hence \eqref{eq:p:1} gives
\[
\begin{aligned}
mW_4
&\ge(m+1)\left(m-\frac{29}{10}\right)-m(m-1)\\
&=-\frac{9m+29}{10}>-m,
\end{aligned}
\]
where the strict inequality is equivalent to \(m>29\). Thus \(W_4>-1\),
and its integrality gives \(W_4\ge0\). More generally, the same calculation
would accept a uniform deficit \(D_0\ge m-C\) for all \(m\ge30\) whenever
\(C<90/31\); the chosen value \(C=29/10\) lies just below this threshold.
Together with the three-plane case and the small-multiplicity argument, this
proves the theorem.
\end{proof}
